\chapter*{致\qquad 谢\markboth{致\qquad 谢}{}}

时光如白驹过隙，转眼间三年的研究生生涯已悄然走到尾声。回首这段求学历程，心中充满了感激与感慨。

首先，衷心感谢我的导师魏显老师在整个研究生阶段对我的悉心指导与无私帮助。导师渊博的学识、严谨的治学态度和对学术的热忱深深影响着我，每当研究遇到困难时，导师的鼓励和建议总能为我指引方向。没有导师的引领，本文的完成是不可能的。

感谢课题组的合作伙伴们。每一次思想的碰撞都让我对双曲几何和路网学习有了更深刻的理解,与你们共同探讨问题、解决困难的经历是我研究生生涯中最珍贵的记忆之一。

感谢华东师范大学软件工程学院给我提供了良好的学习和研究环境。感谢各位任课老师在专业知识上的传授，使我奠定了扎实的理论基础。感谢学院的行政老师们，在生活和学业上给予了诸多帮助。

感谢 VecCity 团队提供的开放式基准框架，为 HyperRoad 的实验评估提供了标准化的实验环境。感谢开放街道地图（OpenStreetMap）社区的所有贡献者，为路网研究提供了宝贵的开源数据基础。


最后，最深切的感谢要献给我的家人。感谢父母多年来无微不至的关爱与支持，是你们的鼓励给了我前行的勇气和力量。无论前路如何，你们永远是我最坚强的后盾。

书山有路勤为径，学海无涯苦作舟。研究生阶段的学习让我深刻认识到科学探索的无尽魅力，也让我对学术研究有了更成熟的思考与理解。本文的工作仅是路网表示学习这一研究领域的一个小小起点，希望在未来的研究道路上能够不断探索，取得更多有意义的成果。

\vspace{0.8cm} \hspace{10cm}  欧阳明杰

\hspace{8cm}  二零二五年三月于上海